\subsection*{الف}
یک \lr{MLP} با یک لایهٔ پنهان و بدون تابع فعال‌سازی، در واقع یک مدل خطی روی ورودی است. اگر
$x \in \mathbb{R}^{d}$،
وزن‌های لایهٔ پنهان و خروجی به‌ترتیب
$W_1 \in \mathbb{R}^{k \times d}$،
$W_2 \in \mathbb{R}^{m \times k}$
باشند، آنگاه
\[
\hat{y} = W_2 (W_1 x) = (W_2 W_1)\,x.
\]
پس تنها چیزی که روی داده اثر می‌گذارد ماتریسِ ضرب‌شدهٔ
\[
M := W_2 W_1 \in \mathbb{R}^{m \times d}
\]
است و مسئلهٔ یادگیری را می‌توان به‌صورتِ بهینه‌سازی روی $M$ دید. برای بسیاری از توابع هزینهٔ رایجِ یادگیری نظارت‌شده که روی \lr{linear predictor} تعریف می‌شوند (مثل کمترین مربعات و \lr{logistic / softmax cross-entropy}) تابع هزینه نسبت به $M$ یک مسئلهٔ کوژ است:
\[
\min_{W_1,W_2} L(W_2W_1) \quad \equiv \quad \min_{M} L(M).
\]

حال نشان می‌دهیم گزارهٔ سؤال نادرست است (یعنی کمینهٔ محلیِ بد نداریم).
فرض کنید $(W_1,W_2)$ یک کمینهٔ محلی باشد و $M=W_2W_1$.
اگر $M$ کمینهٔ سراسریِ مسئلهٔ کوژ روی $M$ نباشد، آنگاه چون $L(M)$ کوژ است، یک جهت کاهشی $\Delta M$ وجود دارد به‌طوری‌که برای $\alpha>0$ کوچک:
\[
L(M+\alpha \Delta M) < L(M).
\]
از طرفی تغییرات کوچکِ $W_1,W_2$ تغییرِ مرتبهٔ اولِ زیر را روی $M$ ایجاد می‌کند:
\[
(W_2+\Delta W_2)(W_1+\Delta W_1) = W_2W_1 + \underbrace{W_2\Delta W_1 + \Delta W_2 W_1}_{\Delta M} + O(\|\Delta W\|^2).
\]
اگر $W_1$ سطر-تمام‌رتبه باشد، می‌توان $\Delta W_1=0$ و
\[
\Delta W_2 := \Delta M\, W_1^{\mathsf T} (W_1 W_1^{\mathsf T})^{-1}
\]
گرفت تا $\Delta W_2 W_1 = \Delta M$ و بنابراین در همسایگیٔ $(W_1,W_2)$ یک جهت کاهشی برای $L$ ساخته شود؛ این با کمینهٔ محلی بودن تناقض دارد.
اگر $W_1$ سطر-تمام‌رتبه نباشد، آنگاه نقطه روی یک منیفلدِ تخت از فاکتورگیری‌ها قرار می‌گیرد (جهت‌هایی وجود دارد که $M$ را تغییر نمی‌دهند) و چنین نقاطی کمینهٔ موضعیِ جداافتاده نیستند؛ در واقع با یک اغتشاش بسیار کوچک می‌توان به فاکتورگیری هم‌ارز (با همان $M$) که سطر-تمام‌رتبه است رسید و سپس استدلال بالا را اعمال کرد. نتیجه:
\[
\text{در این شبکهٔ خطی (بدون فعال‌سازی)، کمینهٔ محلیِ غیرسراسری وجود ندارد و گزاره نادرست است.}
\]

\subsection*{ب}
با وجود غیرمحدب بودن، روش‌های گرادیانی در شبکه‌های عمیق غالباً موفق‌اند چون هندسهٔ نقاط بحرانی در ابعاد بالا به نفعِ فرار از بدترین نواحی است:
\begin{itemize}
    \item در فضاهای پربعد، «بیشتر» نقاط بحرانی از نوع \lr{saddle point} هستند نه کمینهٔ محلیِ بد. یعنی هسین در آن نقاط دارای مقادیر ویژهٔ با علامت‌های مختلف است.
    \item اگر $H=\nabla^2 L(\theta^\star)$ هسین در یک نقطهٔ بحرانی باشد، طیف هسین معمولاً شامل یک \lr{bulk} بزرگ نزدیک به صفر و تعداد کمی \lr{outlier} است؛ بنابراین بسیاری از جهت‌ها «تخت» هستند (مقادیر ویژه نزدیک صفر) و چند جهتِ محدود خمیدگی قوی دارند.
    \item در \lr{saddle}ها وجود مقدار ویژهٔ منفی یعنی یک جهتِ کاهشیِ درجهٔ دوم وجود دارد؛ الگوریتم‌هایی مثل \lr{SGD} به‌خاطر نویز ذاتی مینی‌بچ یا اغتشاشات عددی، با احتمال بالا از این جهت‌ها استفاده کرده و از \lr{saddle} فرار می‌کنند.
    \item کمینه‌های قابل‌دستیابی در شبکه‌های بزرگ اغلب «هم‌ارزش»‌اند (مقدار هزینه نزدیک) و به‌واسطهٔ تقارن‌ها و بیش‌پارامتری بودن، نواحی کم‌هزینه به‌صورت پهن و بعضاً به هم متصل‌اند؛ این باعث می‌شود رسیدن به یک کمینهٔ خوب آسان‌تر از گیر افتادن در کمینهٔ بد باشد.
\end{itemize}

\subsection*{پ}
گرادیان در گام $t$ برابر است با
\[
\nabla L = 
\begin{bmatrix}
100\\
0.01
\end{bmatrix}.
\]

\begin{enumerate}
\item[(۱)]
اگر نرخ یادگیری ثابت $\eta$ باشد، اندازهٔ گام در مؤلفه‌ها به‌ترتیب
\[
\Delta \theta_1 = -100\eta, \qquad \Delta \theta_2 = -0.01\eta
\]
می‌شود. این یک دوگانگی می‌سازد:
\begin{itemize}
    \item اگر $\eta$ را برای پایداری در راستای اول کوچک بگیریم (تا نوسان/واگرایی در راستای با شیب بزرگ رخ ندهد)، در راستای دوم گام بسیار کوچک شده و همگرایی کند می‌شود.
    \item اگر $\eta$ را برای پیشرفت مناسب در راستای دوم بزرگ بگیریم، در راستای اول جهش بسیار بزرگ رخ می‌دهد و ممکن است از مینیمم عبور کنیم یا واگرا شویم.
\end{itemize}
این دقیقاً همان مشکل \lr{ill-conditioning} (ناهم‌مقیاسی مؤلفه‌ها/خمیدگی‌ها) است.

\item[(۲)]
روش‌های تطبیقی مانند \lr{RMSProp} و \lr{Adam} به‌جای یک $\eta$ مشترک، یک مقیاس‌بندی مؤلفه‌به‌مؤلفه اعمال می‌کنند و عملاً گام را بر اساس «بزرگی گرادیان‌های گذشته» نرمال می‌کنند. به‌صورت خلاصه:
\[
v_t = \rho v_{t-1} + (1-\rho)\, g_t \odot g_t
\]
و سپس به‌روزرسانی از فرمِ
\[
\theta_{t+1} = \theta_t - \eta \frac{g_t}{\sqrt{v_t}+\epsilon}
\]
(\lr{RMSProp})
یا در \lr{Adam} با میانگینِ مرتبهٔ اول نیز:
\[
m_t = \beta m_{t-1} + (1-\beta) g_t, \qquad
\theta_{t+1} = \theta_t - \eta \frac{\hat{m}_t}{\sqrt{\hat{v}_t}+\epsilon}.
\]
بنابراین این روش‌ها روی «ریشهٔ میانگین مربعات گرادیان» (برآورد لحظهٔ دوم، یعنی $g^2$) نرمال‌سازی انجام می‌دهند و برای مؤلفهٔ بزرگ (مثل $100$) گام را کوچک‌تر و برای مؤلفهٔ کوچک (مثل $0.01$) گام نسبی را بزرگ‌تر می‌کنند.
\end{enumerate}

\subsection*{ت}
فرض کنید $p \gg n$ و شبکه تا خطای آموزش صفر آموزش داده شود. اکنون $10\%$ برچسب‌ها به‌صورت تصادفی برعکس و شبکه دوباره آموزش داده می‌شود.

\begin{enumerate}
\item[(۱)]
در حالت بیش‌پارامتری‌شده، ظرفیت مدل بسیار بالاست و شبکه می‌تواند داده‌ها را «حفظ» کند. به‌طور شهودی، وقتی پارامترها بسیار بیشتر از قیود ناشی از داده‌هاست، مدل می‌تواند تابعی بسازد که روی همان $n$ نمونه هر الگوی برچسب (حتی الگوی نویزی/تصادفی) را بازتولید کند؛ بنابراین حتی با برچسب‌های غلط نیز می‌توان به خطای آموزش صفر رسید (در عمل نیز شبکه‌های عمیق توانایی \lr{memorization} دارند).

\item[(۲)]
اما روی آزمون، نویز برچسب قابل‌تعمیم نیست. اگر برچسب را به‌شکل
\[
y = y_{\text{signal}} + \epsilon
\]
در نظر بگیریم (که $\epsilon$ نویز ناشی از برچسب‌های برعکس‌شده است)، آموزش تا خطای صفر یعنی مدل ناچار است علاوه بر سیگنال، مؤلفهٔ نویز را هم برازش کند:
\[
f \approx f_{\text{signal}} + f_{\text{noise}}.
\]
جزء $f_{\text{noise}}$ روی دادهٔ آموزش کمک می‌کند خطا صفر شود، ولی چون این نویز با دادهٔ آزمون همبستگی ندارد، باعث افزایش خطای آزمون می‌شود. بیش‌پارامتری‌بودن این اثر را تشدید می‌کند چون مدل برای برازش نویز نیز آزادی کافی دارد (مرز تصمیم/تابع تقریب پیچیده‌تر می‌شود)؛ در نتیجه با وجود دقت کامل روی آموزش، تعمیم‌پذیری افت می‌کند.
\end{enumerate}